\section{An exact dominant and alternating sector decomposition}\label{tks:sec:sectors}
The fixed-order expansion can be refined without claiming a complete transseries or resurgence theory. The normal-moment representation of involution numbers admits an exact split over the two half axes. Each half has its own fixed-order asymptotic expansion, and the global support method transfers those separate expansions to the Kostka sum.

\subsection{Exact sector definitions}
For $\sigma\in\{+1,-1\}$, define
\begin{equation}\label{tks:eq:halfaxis}
\mathcal M_m^\sigma=\frac{e^{-1/2}}{\sqrt{2\pi}}
 \int_0^\infty x^m e^{-x^2/2+\sigma x}\,\dd x,
\qquad m\geq0.
\end{equation}
If $X$ is normal with mean $1$ and variance $1$, its moment generating function is $e^{z+z^2/2}$. Thus
\begin{equation}\label{tks:eq:exactIsplit}
I_m=\E X^m=\mathcal M_m^++(-1)^m\mathcal M_m^-.
\end{equation}
Write
\[
w_{n,s}=\sum_{\substack{\rho\vdash s\\\min\rho\geq2}}
                \alpha_{1^{n-s}\rho}r(\rho),
\qquad
\mathcal A_n^\sigma=\sum_{s=0}^n\sigma^s w_{n,s}\mathcal M_{n-s}^\sigma.
\]
The empty type is included. Each $w_{n,s}$ is nonnegative, but $\mathcal A_n^-$ is defined by a signed sum. Inserting \eqref{tks:eq:exactIsplit} in \eqref{tks:eq:schwob} gives the exact identity
\begin{equation}\label{tks:eq:exactasplit}
a_n=\mathcal A_n^++(-1)^n\mathcal A_n^-\qquad(n\geq0).
\end{equation}
The two sectors are canonical relative to this specified half-axis split; no positivity of every recessive-sector summand is asserted.

Let
\[
B_\sigma(n)=\frac1{\sqrt2}\left(\frac ne\right)^{n/2}
                          e^{\sigma\sqrt n-1/4},
\qquad
\mathcal K(t)=\sum_{j\geq0}k_jt^j
\]
be the formal full Kostka series normalized by $CB_+(n)$. Its coefficients through $k_6$ are those in \eqref{tks:eq:amultiplier}; in particular $k_0=1$, $k_1=7/24$, and $k_2=-119/1152$.

\begin{theorem}\label{tks:thm:sectors}
For every fixed $S\geq0$ and either sign,
\begin{equation}\label{tks:eq:sectorsupport}
\mathcal A_n^\sigma
 =C\sum_{s=0}^{S}\sigma^sH_s\mathcal M_{n-s}^\sigma
  +O_S\bigl(\mathcal M_n^\sigma n^{-(S+1)/2}\bigr).
\end{equation}
For every fixed $K\geq0$, the separate sector expansions are
\begin{equation}\label{tks:eq:sectorseries}
\mathcal A_n^\sigma=CB_\sigma(n)
 \left(\sum_{j=0}^{K}k_j\sigma^j n^{-j/2}
                    +O_K(n^{-(K+1)/2})\right).
\end{equation}
In particular, $\mathcal A_n^-$ is eventually positive, and
\begin{equation}\label{tks:eq:sectorratio}
\frac{\mathcal A_n^-}{\mathcal A_n^+}
 =e^{-2\sqrt n}\left(1-\frac7{12\sqrt n}
                         +\frac{49}{288n}+O(n^{-3/2})\right).
\end{equation}
\end{theorem}

The exact half-axis representation is a specialization of classical parabolic-cylinder integrals. In the notation of DLMF \cite[(12.5.1)]{DLMF},
\begin{equation}\label{tks:eq:parabolic}
\mathcal M_n^\sigma
  =\frac{n!e^{-1/4}}{\sqrt{2\pi}}U(n+\tfrac12,-\sigma).
\end{equation}
Olver's separate large-order expansions cover these two branches \cite[\S11]{Olver}. The purpose of the elementary proof below is the transfer through the Kostka support weights, with precise normalization and remainder scope; it is not a novelty claim for the classical involution sectors.

\subsection{Uniform half axis moment ratios}
Let $R_k^\sigma=\mathcal M_k^\sigma/\mathcal M_{k-1}^\sigma$ for $k\geq1$. Moment log-convexity from Cauchy--Schwarz and integration by parts give
\begin{equation}\label{tks:eq:momentrec}
R_{k+1}^\sigma\geq R_k^\sigma,\qquad
R_{k+1}^\sigma=\sigma+\frac{k}{R_k^\sigma}\quad(k\geq1).
\end{equation}
For the minus sign, these imply $(R_k^-)^2+R_k^-\leq k$, hence $R_k^-\leq\sqrt k$. For $k\geq2$ the recurrence then yields
\[
R_k^-\geq\sqrt{k-1}-1\geq\tfrac12\sqrt k\quad(k\geq9).
\]
The finitely many earlier ratios are strictly positive. For the plus sign, $(R_k^+)^2-R_k^+\leq k$ gives $R_k^+\leq1+\sqrt k$, and hence
\[
R_k^+\geq1+\frac{k-1}{1+\sqrt{k-1}}\geq\sqrt{k-1}
\quad(k\geq2).
\]
Together with the positive first ratios, there is an absolute $c_*>0$ such that $R_k^\sigma\geq c_*\sqrt k$ for both signs and all $k\geq1$. Therefore
\begin{equation}\label{tks:eq:momentbound}
\frac{\mathcal M_{n-s}^\sigma}{\mathcal M_n^\sigma}
 \leq c_*^{-s}\fall ns^{-1/2}
 \leq\left(\frac{\sqrt e}{c_*\sqrt n}\right)^s.
\end{equation}
Using \eqref{tks:eq:Bns} and \eqref{tks:eq:rootbound}, even the absolute signed-sector support contribution is bounded by
\[
w_{n,s}\frac{\mathcal M_{n-s}^\sigma}{\mathcal M_n^\sigma}
 \leq C_0\left(4\sqrt{2e}\,c_*^{-1}n^{-1/4}\right)^s.
\]
The tail from $T$ is $O_T(n^{-T/4})$. Fixed-support stabilization still gives $w_{n,s}=CH_s+O_{s,R}(R^{-n})$. Splitting omitted supports at $T=2S+2$ proves \eqref{tks:eq:sectorsupport}, exactly as in Section~\ref{tks:sec:support}, without relying on cancellation.

\subsection{A real saddle and coefficient parity}
Set $N=\sqrt n$, $t=1/N$, and substitute $x=N+y$ in \eqref{tks:eq:halfaxis}. Exact simplification gives
\begin{align}
\frac{\mathcal M_n^\sigma}{B_\sigma(n)}
 &=\frac1{\sqrt\pi}\int_{-N}^{\infty}
     e^{-(y-\sigma/2)^2}e^{\Psi_N(y)}\,\dd y,\label{tks:eq:realmoment}\\
\Psi_N(y)&=N^2\log(1+y/N)-Ny+y^2/2.\notag
\end{align}
Locally, for $|ty|<1$,
\[
\Psi_N(y)=\sum_{m\geq1}t^m v_m(y),\qquad
v_m(y)=\frac{(-1)^{m+1}y^{m+2}}{m+2}.
\]
Let $V_0(y)=1$ and $V_j(y)=j^{-1}\sum_{m=1}^jmv_m(y)V_{j-m}(y)$. Then the half-axis coefficients are the rational Gaussian moments
\begin{equation}\label{tks:eq:sectorcoeff}
\alpha_j^\sigma=\frac1{\sqrt\pi}
       \int_\R e^{-(y-\sigma/2)^2}V_j(y)\,\dd y.
\end{equation}
Every $v_m$ has parity $m$, so $V_j(-y)=(-1)^jV_j(y)$. Reflection gives
\begin{equation}\label{tks:eq:sectorparity}
\alpha_j^-=(-1)^j\alpha_j^+.
\end{equation}
The plus Gaussian moments are computed by replacing $y^d$ with
\[
\sum_{r=0}^{\lfloor d/2\rfloor}
  \binom d{2r}(1/2)^{d-2r}\frac{(2r-1)!!}{2^r},
\]
so this is again a finite rational algorithm at each fixed order.

For completeness, the local expansion has a controlled remainder. Put
$f_N(y)=N^2\log(1+y/N)-Ny-y^2/2$, for $y>-N$. Then
$f_N(0)=f_N'(0)=0$ and
\[
f_N''(y)=-\frac{N^2}{(N+y)^2}-1\leq-1,
\qquad f_N(y)\leq-y^2/2.
\]
Consequently the exact normalized integrand outside $|y|\leq N^\eps$ is bounded by a Gaussian tail after absorbing the linear term $\sigma y$, for fixed $0<\eps<1/3$. On the central interval, Taylor's theorem through order $K$ bounds the remainder by $t^{K+1}$ times a fixed polynomial in $|y|$ times $e^{\delta y^2}$, for arbitrarily small fixed $\delta>0$ once $N$ is large. Here $|ty|=o(1)$, fixed derivatives are polynomially bounded, and the perturbation is $O(t|y|^3)\leq\delta y^2$. Gaussian integration and extension of the polynomial integrals to $\R$ give
\begin{equation}\label{tks:eq:momentexpansion}
\mathcal M_n^\sigma=B_\sigma(n)
 \left(\sum_{j=0}^{K}\alpha_j^\sigma n^{-j/2}
                 +O_K(n^{-(K+1)/2})\right).
\end{equation}
The leading terms show $\mathcal M_n^-/\mathcal M_n^+\sim e^{-2\sqrt n}$. By \eqref{tks:eq:exactIsplit} and uniqueness of the dominant Poincar\'e expansion, $\alpha_j^+$ are exactly the coefficients $\alpha_j$ in \eqref{tks:eq:Ialpha}.

\subsection{Composition and the separate remainders}
Let $\mathcal I(t)=\sum_{j\geq0}\alpha_jt^j$, and define
\[
E_s(t)=\frac{t^{-2}-s}{2}\log(1-st^2)+\frac s2
                        +\frac{\sqrt{1-st^2}-1}{t}.
\]
The apparent singularities are removable. For the plus sector, the coefficient rule is
\begin{equation}\label{tks:eq:sectorcomposition}
\mathcal K(t)=\sum_{s\geq0}H_st^s e^{E_s(t)}
                    \mathcal I\left(\frac{t}{\sqrt{1-st^2}}\right).
\end{equation}
Only $s\leq K$ enters a coefficient of degree $K$, so this is finite at every requested order. In the minus sector, the signed support weight $(-1)^s$, the reversed sign in the square-root exponential, and \eqref{tks:eq:sectorparity} combine to replace $t$ by $-t$ in \eqref{tks:eq:sectorcomposition}. Using \eqref{tks:eq:sectorsupport} at $S=K$ and the finitely many shifted versions of \eqref{tks:eq:momentexpansion} proves \eqref{tks:eq:sectorseries}, with its separate relative error for each sector. Dividing the two expansions gives \eqref{tks:eq:sectorratio}, since
\[
\frac{\mathcal K(-t)}{\mathcal K(t)}
 =1-2k_1t+2k_1^2t^2+O(t^3).
\]

This proves an exact alternating contribution $(-1)^n\mathcal A_n^-$ with its own all-fixed-order expansion. A fixed finite dominant truncation provides only a remainder bound of order $B_+(n)n^{-(K+1)/2}$, which is larger than $B_-(n)$. That available bound does not isolate the recessive term by subtraction from $a_n$; it is not a lower bound on every actual truncation error. The theorem neither classifies all smaller effects inside the exact sectors nor establishes optimal truncation, Stokes multipliers, or global resurgence.
\begin{writenote}[The two sectors at small $n$, 7 October 2026]
For $15\le n\le39$ the write compared $a_n$ (b-file) with $CB_+(n)\mathcal K(t)$
truncated after $t^6$, with the constants above. Without the recessive sector
the remainder, times $n^{7/2}$, alternates with the parity of~$n$ ($-14.3$,
$-17.3$ at $n=38$, $39$); adding
$(-1)^nCB_-(n)\mathcal K(-t)$, as~\eqref{tks:eq:exactasplit}
and~\eqref{tks:eq:sectorseries} prescribe, removes the alternation ($-15.7$,
$-15.9$). The alternating sector, of relative size $e^{-2\sqrt n}$, is thus
already visible in the exact values; this is an observation, not a
certificate.
[Independent check, 7~October 2026. The numbers were reproduced, and a
precision is needed. With the recessive sector added, the scaled remainder
increases monotonically for $15\le n\le34$. For $34\le n\le39$ a residual
alternation of about $\pm0.1$ remains ($-15.79$, $-15.87$, $-15.61$, $-15.79$,
$-15.75$, $-15.98$), down from about $\pm1.7$.

The residual comes from a third exponentially small effect, which both
expansions omit because it lies beyond all their orders. The support weight
$\alpha_{1^n}=[x^n]\prod_{j\ge1}(1-x^j/j!)^{-1}$ tends to~$C$ with error
$O(2^{-n/2})$, from the poles $x=\pm\sqrt2$ of the factor $j=2$. The pole at
$-\sqrt2$ gives $a_n$ an alternating relative term of about
$0.0644\,(-1)^n2^{-n/2}$, comparable with $e^{-2\sqrt n}$ near $n=33$ and
larger below.

In units of $n^{-7/2}$, the exact correction $(\alpha_{1^n}-C)I_n/(CB_+(n))$ is
$-92.8$ at $n=15$ and $-0.63$ at $n=39$. Subtracting it as well leaves a remainder
that rises from $-51.9$ at $n=15$ to about $-13.0$ at $n=28$--$30$ and falls to
$-15.3$ at $n=39$. Its only parity wiggle is one of about $0.1$ at the turning
point, $n=28$--$30$.]
\end{writenote}
